\documentclass[11pt,a4paper]{article}

\usepackage[utf8]{inputenc}
\usepackage[T1]{fontenc}
\usepackage[english]{babel}
\usepackage{amsmath,amssymb,amsthm}
\usepackage[margin=2.5cm]{geometry}
\usepackage{booktabs}
\usepackage{hyperref}
\usepackage{xcolor}
\hypersetup{colorlinks=true, linkcolor=blue!50!black, citecolor=blue!50!black, urlcolor=blue!50!black}

\newcommand{\RV}{\mathrm{RV}}
\newcommand{\IV}{\mathrm{IV}}
\newcommand{\E}{\mathbb{E}}

\title{Volatility Forecasting with Machine Learning\\
and Trading the Volatility Risk Premium\\
\large Theoretical foundations of the \texttt{vol\_ml\_fund} project}
\author{Quantitative research project}
\date{\today}

\begin{document}

\maketitle

\begin{abstract}
This document presents the theoretical foundations of the
\texttt{vol\_ml\_fund} research platform: the measurement of realized
volatility and its estimators, implied volatility and the variance risk
premium, forecasting models (HAR-RV and machine learning), the purged
walk-forward evaluation protocol, and the transformation of forecasts into
trading positions. Each section justifies the design choices made in the code
and, in particular, the improvements introduced in version~2 of the project.
This document is educational: it constitutes neither investment advice nor a
promise of performance.
\end{abstract}

\tableofcontents

%=============================================================================
\section{Introduction: why trade volatility?}

Volatility exhibits two statistical properties that set it apart from returns
themselves: it is \emph{persistent} (volatility clustering: turbulent periods
follow turbulent periods) and partly \emph{predictable}, whereas the sign of
returns is essentially not. This is what makes volatility forecasting
attractive for a machine learning project: the signal-to-noise ratio is far
more favorable than for directional prediction.

The economic foundation of the strategy studied here is the \emph{volatility
risk premium} (or variance risk premium): on average, the volatility implied
by option prices exceeds the volatility that is subsequently realized.
Sellers of crash insurance get paid a premium. The project consists of
(i)~forecasting future realized volatility, (ii)~comparing it with the
implied volatility quoted by the market, and (iii)~taking a position when the
predicted gap deviates sufficiently from its usual value.

%=============================================================================
\section{Measuring realized volatility}

\subsection{Logarithmic returns}

For a closing price $P_t$, the log return is
\begin{equation}
r_t = \ln\frac{P_t}{P_{t-1}}.
\end{equation}
Log returns are additive across time
($r_{t,t+k} = \sum_{i=1}^{k} r_{t+i}$), which simplifies multi-horizon
aggregation, and they fit naturally with the assumption of log-normal prices.
\emph{Implementation:} \texttt{src/data/preprocess.py}.

\subsection{Close-to-close estimator}

The annualized realized volatility over an $n$-day window is
\begin{equation}
\RV_t^{(n)} = \sqrt{\frac{252}{n-1}\sum_{i=0}^{n-1}\bigl(r_{t-i}-\bar{r}\bigr)^2},
\end{equation}
i.e.\ the rolling standard deviation of daily returns multiplied by
$\sqrt{252}$ (the number of trading days per year). The window ends at date
$t$: no future information is used.

\subsection{Range-based estimators: Parkinson and Garman--Klass}
\label{sec:range}

The close-to-close estimator uses only one data point per day and discards
the intraday information contained in the daily extremes. Parkinson (1980)
shows that with the daily high $H_t$ and low $L_t$,
\begin{equation}
\hat{\sigma}^2_{P,t} = \frac{1}{4\ln 2}\,\Bigl(\ln\frac{H_t}{L_t}\Bigr)^2
\end{equation}
is roughly \textbf{5 times more efficient} (its sampling variance is 5 times
smaller) than $r_t^2$, under driftless Brownian motion. Garman and Klass
(1980) add the open $O_t$ and the close $C_t$:
\begin{equation}
\hat{\sigma}^2_{GK,t} = \tfrac{1}{2}\Bigl(\ln\tfrac{H_t}{L_t}\Bigr)^2
- (2\ln 2 - 1)\Bigl(\ln\tfrac{C_t}{O_t}\Bigr)^2,
\end{equation}
gaining further efficiency ($\approx 7.4\times$). Annualization is obtained
by averaging over the window and multiplying by 252 before taking the square
root.

\textbf{Why this improvement:} at daily frequency, this is the cheapest way
to increase the precision of the volatility features --- the statistical
equivalent of several extra days of observations, without high-frequency
data. \emph{Implementation:} \texttt{garman\_klass\_volatility} and
\texttt{parkinson\_volatility} in
\texttt{src/features/volatility\_features.py}.

\subsection{Downside semi-volatility and the leverage effect}

Volatility reacts asymmetrically to returns: a price decline increases future
volatility more than a rise of the same magnitude (the \emph{leverage
effect}). The downside semi-volatility
\begin{equation}
\RV^{-}_t = \sqrt{\frac{252}{n}\sum_{i=0}^{n-1} r_{t-i}^2\,
\mathbf{1}\{r_{t-i}<0\}}
\end{equation}
captures this asymmetry; adding it to the regressors follows the spirit of
the LHAR model (HAR with leverage) and of Patton and Sheppard (2015) on
``good'' and ``bad'' volatility.

%=============================================================================
\section{Implied volatility: the VIX and its term structure}

The VIX measures the 30-day implied volatility of S\&P~500 options, computed
model-free from a continuum of strikes (CBOE methodology; Whaley, 2009).
Divided by 100, it is directly comparable to an annualized realized
volatility. The VIX3M is its 3-month counterpart.

The ratio $\mathrm{VIX}/\mathrm{VIX3M}$ summarizes the \emph{term structure}
of implied volatility:
\begin{itemize}
\item ratio $<1$ (\emph{contango}): calm regime, short-dated vol below
  long-dated vol --- the normal state, favorable to \emph{short vol}
  strategies;
\item ratio $>1$ (\emph{backwardation}): market stress, demand for immediate
  hedging --- often heralds elevated realized volatility.
\end{itemize}
It is one of the most extensively documented predictors of volatility
regimes, hence its inclusion among the features
(\texttt{vix\_term\_structure}).

%=============================================================================
\section{The variance risk premium}
\label{sec:vrp}

\subsection{Definition and stylized facts}

The variance risk premium (VRP) is the gap between implied variance and the
expectation of future realized variance:
\begin{equation}
\mathrm{VRP}_t = \IV_t^2 - \E_t\!\left[\RV^2_{t\to t+h}\right] .
\end{equation}
Empirically it is \emph{positive on average}: on equity indices, $\IV$
exceeds future RV roughly 85\,\% of the time (Carr and Wu, 2009; Bollerslev,
Tauchen and Zhou, 2009). Systematically selling volatility collects this
premium --- at the price of rare but brutal losses when volatility explodes
(the ``picking up pennies in front of a steamroller'' profile).

\subsection{Why target the ratio rather than the level}
\label{sec:target}

Version~1 of the project predicted the \emph{level} of future RV. The Random
Forest diagnostic was unambiguous: implied volatility carried 90\,\% of the
feature importance. The model was essentially learning that
``$\RV_{\text{future}} \approx \mathrm{VIX}$'', information already known to
the entire market and therefore without trading value.

Version~2 predicts instead
\begin{equation}
y_t = \ln\frac{\RV_{t+1\to t+h}}{\IV_t},
\label{eq:target}
\end{equation}
for three reasons:
\begin{enumerate}
\item \textbf{It is the traded quantity.} The signal compares precisely the
  predicted RV with the quoted IV; predicting their log-ratio directly aligns
  the model's loss function with the economic object of interest (the VRP
  residual: what the option market does \emph{not} already price).
\item \textbf{The logarithm stabilizes the distribution.} RV is positive and
  strongly right-skewed (fat tail: March 2020, August 2015\ldots). In log
  space the distribution is close to normal (Andersen et al., 2001), the
  error variance is more homogeneous, and extreme episodes no longer dominate
  least-squares estimation.
\item \textbf{The ratio is approximately stationary.} The level of volatility
  switches regimes (5\,\% in 2017, 80\,\% in 2020); the RV/IV ratio
  oscillates around a stable mean ($\approx e^{-0.2}$), which helps
  out-of-sample generalization.
\end{enumerate}
Predictions are mapped back to levels via
$\widehat{\RV} = \IV_t\, e^{\hat{y}}$ for evaluation, so that all models and
benchmarks are compared on the same scale.
\emph{Implementation:} \texttt{transform\_target} / \texttt{invert\_target}.

%=============================================================================
\section{Forecasting models}

\subsection{HAR-RV: the reference benchmark}

The HAR-RV model (\emph{Heterogeneous AutoRegressive}, Corsi 2009) regresses
future volatility on realized volatilities measured over three horizons:
\begin{equation}
\RV_{t+h} = \beta_0 + \beta_d\, \RV_t^{(d)} + \beta_w\, \RV_t^{(w)}
+ \beta_m\, \RV_t^{(m)} + \varepsilon_{t+h},
\end{equation}
where $d$, $w$, $m$ denote daily, weekly and monthly windows (here
approximated by 5, 20 and 60 days in the absence of intraday data). The
intuition is that of a heterogeneous \emph{volatility cascade}: traders
operate at different horizons and long-horizon volatility influences
short-horizon volatility. Despite its simplicity (three OLS coefficients),
the HAR is notoriously hard to beat --- which is why it serves as the
reference in the Diebold--Mariano tests.

\subsection{HAR-X: making the benchmark honest}

Comparing a HAR that sees no option-market information against ML models that
see the VIX would be unfair. The \textbf{HAR-X} adds implied volatility as an
exogenous regressor: it is the standard academic benchmark whenever option
data is available, and the real opponent to beat before machine learning can
be declared useful.

\subsection{Regularized regression: Ridge and Lasso}

With some twenty correlated features (the multi-window RVs resemble each
other), OLS becomes unstable. Regularization adds a penalty on the norm of
the coefficients:
\begin{equation}
\hat{\beta} = \arg\min_\beta\; \|y - X\beta\|_2^2
+ \lambda \|\beta\|_2^2 \quad (\text{Ridge}), \qquad
\lambda \|\beta\|_1 \quad (\text{Lasso}).
\end{equation}
Ridge shrinks correlated coefficients towards common values; Lasso sets some
of them to zero (variable selection). In volatility forecasting, these
regularized linear models are frequently on par with ensemble methods at a
fraction of the computational cost, and they guard against the hasty
conclusion that ``non-linear ML is necessary''. Features are standardized
inside a pipeline so that the penalty applies on comparable scales.

\subsection{Random forests and gradient boosting}

The two ensemble methods capture non-linearities (threshold effects on the
term structure, drawdown $\times$ volatility interactions) without manual
specification:
\begin{itemize}
\item \textbf{Random Forest}: an average of many trees decorrelated through
  bootstrapping and feature subsampling; it reduces variance and is robust to
  hyperparameters, but cannot extrapolate beyond the values seen in training
  (an unexpected advantage in log-ratio space, where it bounds extreme
  forecasts).
\item \textbf{Gradient Boosting}: sequential addition of shallow trees
  correcting the residuals of the current model; lower bias, but more prone
  to overfitting --- hence a low learning rate and limited depth.
\end{itemize}
Tree depth and the minimum number of observations per leaf are deliberately
restrictive: in financial time series the signal-to-noise ratio is low and
complexity is paid for out of sample.

\subsection{Quantile regression: quantifying uncertainty}

Gradient boosting with the \emph{pinball} loss estimates conditional
quantiles $q_{10}$ and $q_{90}$ of the target. The band $[q_{10}, q_{90}]$
measures forecast uncertainty: when this band straddles zero (in log-ratio
space), the \emph{sign} of the edge is uncertain and the position size is
halved (section~\ref{sec:signal}). Quantiles in log-ratio space carry over to
RV space through the monotonic transform $x \mapsto \IV e^{x}$, which
preserves quantiles.

\subsection{Ensemble and naive benchmarks}

The \textbf{ensemble} averages the predictions (in RV space) of several model
families: since errors from different families are imperfectly correlated,
averaging reduces variance at essentially no cost in bias --- a classic
result of the forecast combination literature (Bates and Granger, 1969).

Two \textbf{naive benchmarks} are always included:
$\widehat{\RV} = \RV_t^{(5)}$ (random walk in volatility) and
$\widehat{\RV} = \IV_t$ (the option market is right). A model unable to beat
these two does not deserve further attention: this is the fundamental
non-regression test of any forecasting project.

%=============================================================================
\section{Evaluation protocol}

\subsection{Look-ahead bias}

A backtest is only meaningful if, at every date, the strategy uses only the
information available on that date. Three classic leaks are blocked and
\emph{unit-tested} (\texttt{tests/test\_lookahead.py}):
\begin{enumerate}
\item \textbf{Target construction}: $y_t$ uses returns from $t+1$ to $t+h$
  exclusively; features stop at $t$.
\item \textbf{Execution}: a position decided at the close of $t$ only touches
  the return of $t+1$ (one-day lag in the backtest engine).
\item \textbf{Target overlap} (purge): see below.
\end{enumerate}

\subsection{Walk-forward with purging}
\label{sec:wf}

A single train/test split (v1) suffers from two flaws: the model is frozen
for years (nobody does that in practice), and the performance measure depends
heavily on the chosen cut-off. The \emph{walk-forward} protocol re-fits every
model on a rolling window (1260 days $\approx$ 5 years) every 63 days
($\approx$ one quarter), and predicts only the next block.

\textbf{Purging} is the subtle point: the target at date $t$ contains returns
up to $t+h$. If training runs up to $t$ and prediction starts at $t+1$, the
last $h$ training targets overlap the predicted period --- a genuine, if
discreet, information leak. We therefore leave a gap of $h$ days between the
end of the training set and the first prediction (the general
purging/embargo rule popularized by López de Prado, 2018).

\subsection{Metrics and sub-period stability}

Forecasts are evaluated in RV space with RMSE, MAE and $R^2$. The RMSE is
also computed \textbf{per calendar year}: a model may post a good global
score merely because it navigated (or failed) a single episode --- 2020
typically dominates any global metric. The annual table makes this phenomenon
immediately visible.

\subsection{The Diebold--Mariano test}

Comparing two RMSEs says nothing about whether the gap is significant. The
Diebold--Mariano test (1995) considers the loss differential
$d_t = e_{A,t}^2 - e_{B,t}^2$ and tests $H_0 : \E[d_t]=0$ with the statistic
\begin{equation}
\mathrm{DM} = \frac{\bar{d}}{\sqrt{\widehat{\mathrm{Var}}_{LR}(\bar d)}}
\;\xrightarrow{d}\; \mathcal{N}(0,1),
\end{equation}
where the long-run variance is estimated with Bartlett weights up to lag
$h-1$ (Newey--West): for an $h$-day horizon, forecast errors overlap
mechanically and are autocorrelated up to order $h-1$. A negative and
significant DM against the HAR justifies a model's extra complexity; an
insignificant DM is an equally valid research finding.

%=============================================================================
\section{From signal to portfolio}
\label{sec:signal}

\subsection{The score}

The signal rests on the score
\begin{equation}
s_t = \ln\frac{\widehat{\RV}_{t+1\to t+h}}{\IV_t},
\end{equation}
directly inherited from target~\eqref{eq:target}: $s_t>0$ means the model
expects more volatility than the market prices (volatility is ``cheap''
$\Rightarrow$ long position); $s_t<0$, the opposite. Since the VRP is
positive on average, $s_t$ is negative most of the time: the strategy spends
most of its life \emph{short vol}, collecting the premium, but can flip when
the model anticipates a spike.

\subsection{Hysteresis: two thresholds beat one}

A single threshold (v1) opens and closes the position every time the score
oscillates around it --- with 40\,\% daily turnover, transaction costs
devoured the P\&L. \textbf{Hysteresis} separates the entry threshold from the
exit threshold: a position opens when $|s_t|$ exceeds
$\theta_{\text{entry}}$ and only closes when $|s_t|$ falls back below
$\theta_{\text{exit}} < \theta_{\text{entry}}$. The asymmetric dead band
eliminates the round-trips around the threshold (observed turnover drops from
0.40 to 0.07 per day in our runs).

\subsection{Proportional sizing and conviction}

Rather than a binary $\pm1$ position, size grows with conviction:
\begin{equation}
w_t = d_t \cdot \min\!\Bigl(\frac{|s_t|}{\kappa},\, 1\Bigr)
\cdot \underbrace{\bigl(\tfrac{1}{2}\bigr)^{\mathbf{1}\{q_{10,t}<0<q_{90,t}\}}}_{\text{quantile uncertainty}},
\end{equation}
where $d_t \in \{-1,0,1\}$ is the hysteresis direction and $\kappa$ the
sizing scale. A small edge takes a small position; an edge whose very sign is
uncertain (quantile band straddling zero) is halved.

\subsection{Volatility targeting}

Volatility proxies are themselves enormously and variably volatile
(50--100\,\% annualized). Without control, portfolio risk would be dictated
by the proxy's regime rather than by the model's conviction.
\emph{Volatility targeting} normalizes:
\begin{equation}
w_t^{\text{final}} = w_t \cdot
\min\!\Bigl(\frac{\sigma^\star}{\hat{\sigma}_{\text{proxy},t}},\; L_{\max}\Bigr),
\end{equation}
with $\sigma^\star$ the target volatility (30\,\% in the config),
$\hat{\sigma}_{\text{proxy},t}$ the rolling volatility of the traded leg
(computed from past data only) and $L_{\max}$ the maximum leverage.

\subsection{Instruments: why VIXY and SVXY, and their dangers}

The \emph{spot} VIX is not tradable: it is a calculated index with no cash
replicating portfolio. Version~1 backtested on its returns, producing results
with no economic meaning. Version~2 uses two genuinely tradable ETPs:
\begin{itemize}
\item \textbf{VIXY} (long vol): replicates a rolling basket of
  $\sim$1-month VIX futures. In contango (85\,\% of the time), the
  systematic roll \emph{costs} money: the ETP structurally decays ($-40$ to
  $-70\,\%$ per year in calm regimes). That is the price of insurance.
\item \textbf{SVXY} (short vol): inverse exposure ($-0.5\times$ since 2018)
  to the same basket; it \emph{collects} the roll in contango. A long
  position in SVXY expresses volatility selling without short selling.
\end{itemize}
On February 5, 2018 (the \emph{volmageddon}), the VIX more than doubled in a
single session; XIV (a product analogous to SVXY) lost $\sim$95\,\% and was
liquidated, and SVXY cut its leverage from $-1$ to $-0.5$. This episode is
\emph{inside} our backtest data --- deliberately: any short-vol strategy must
be judged through it. It illustrates the fundamental limitation of the Sharpe
ratio for strategies with fat, asymmetric tails (next section).

\subsection{Transaction costs}

Every weight change on every leg pays a proportional cost $c\,|\Delta w_t|$
with $c = 10$ basis points --- the order of magnitude of the spread plus
impact for liquid ETPs at research-size trades. This model ignores financing,
securities borrowing and the ETPs' management fees (already partly reflected
in their prices).

%=============================================================================
\section{Performance metrics}

For a series of daily net returns $R_t$ and an equity curve
$E_t = E_0\prod_{u\le t}(1+R_u)$:
\begin{align}
\text{Sharpe} &= \frac{\bar{R} - r_f/252}{\sigma_R}\sqrt{252},\\[2pt]
\text{Max drawdown} &= \min_t \Bigl(\frac{E_t}{\max_{u\le t}E_u} - 1\Bigr),\\[2pt]
\text{CAGR} &= \Bigl(\frac{E_T}{E_0}\Bigr)^{252/T} - 1 .
\end{align}
A caveat: for a short-vol strategy, the return distribution is strongly
left-skewed. The Sharpe ratio, which uses only the first two moments,
overstates the quality of such strategies; the max drawdown and a visual
inspection of the equity curve around stress episodes (February 2018,
March 2020) are indispensable to any judgment.

%=============================================================================
\section{Summary: rationale for the v2 improvements}

\begin{center}
\small
\begin{tabular}{p{4.2cm} p{6.2cm} p{3.8cm}}
\toprule
\textbf{Improvement} & \textbf{Theoretical rationale} & \textbf{Module} \\
\midrule
Target $\ln(\RV/\IV)$ & Stationarity, log-normality, alignment with the traded quantity (\S\ref{sec:target}) & \texttt{features/} \\
Garman--Klass, Parkinson & $5$--$7\times$ statistical efficiency (\S\ref{sec:range}) & \texttt{features/} \\
Downside semi-vol & Leverage effect, LHAR & \texttt{features/} \\
VIX/VIX3M term structure & Documented regime predictor & \texttt{features/} \\
Shocks, drawdown, seasonality & Volatility memory, asymmetry & \texttt{features/} \\
HAR-X, Ridge, Lasso & Honest benchmarks; regularization against collinearity & \texttt{models/} \\
Quantile GB & Forecast uncertainty $\to$ sizing & \texttt{models/} \\
Ensemble + naive benchmarks & Forecast combination; non-regression test & \texttt{models/} \\
Purged walk-forward & Realistic re-training; target-overlap leakage (\S\ref{sec:wf}) & \texttt{models/walkforward} \\
Annual RMSE + Diebold--Mariano & Stability across regimes; statistical significance & \texttt{models/evaluation} \\
Hysteresis & Turnover $\div 6$, costs under control & \texttt{signals/} \\
Conviction sizing + vol targeting & Risk driven by the model, not by the proxy & \texttt{signals/} \\
VIXY/SVXY legs & Genuinely tradable instruments, explicit roll yield & \texttt{backtest/} \\
Anti-look-ahead tests & The three classic leaks locked down by pytest & \texttt{tests/} \\
\bottomrule
\end{tabular}
\end{center}

%=============================================================================
\section{Limitations and extensions}

This framework remains a deliberate simplification. The main limitations, in
order of importance:
\begin{enumerate}
\item \textbf{Instruments}: ETPs aggregate the futures roll; working directly
  on the VIX futures curve (or delta-hedged straddles) would give clean
  control over which point of the term structure is bought or sold.
\item \textbf{Hyperparameters}: they are fixed in the config; selecting them
  by temporal cross-validation \emph{nested inside each walk-forward window}
  would be more rigorous (but costly).
\item \textbf{Data snooping}: the signal thresholds and windows were chosen
  reasonably but without a formal procedure; a sensitivity analysis
  (threshold grids, White's reality check) is required before drawing any
  conclusion.
\item \textbf{Tail risk}: no stop-loss, no daily loss limit, no CVaR-based
  sizing --- all indispensable for short vol beyond the research stage.
\item \textbf{Microstructure}: constant costs, no volatility-dependent
  slippage (yet it is precisely when volatility explodes that spreads
  widen).
\end{enumerate}

%=============================================================================
\begin{thebibliography}{99}

\bibitem{andersen2001} Andersen, T., Bollerslev, T., Diebold, F., Labys, P.
(2001). \emph{The Distribution of Realized Exchange Rate Volatility}. Journal
of the American Statistical Association, 96(453).

\bibitem{bates1969} Bates, J., Granger, C. (1969). \emph{The Combination of
Forecasts}. Operational Research Quarterly, 20(4).

\bibitem{btz2009} Bollerslev, T., Tauchen, G., Zhou, H. (2009). \emph{Expected
Stock Returns and Variance Risk Premia}. Review of Financial Studies, 22(11).

\bibitem{carrwu2009} Carr, P., Wu, L. (2009). \emph{Variance Risk Premiums}.
Review of Financial Studies, 22(3).

\bibitem{corsi2009} Corsi, F. (2009). \emph{A Simple Approximate Long-Memory
Model of Realized Volatility}. Journal of Financial Econometrics, 7(2).

\bibitem{dm1995} Diebold, F., Mariano, R. (1995). \emph{Comparing Predictive
Accuracy}. Journal of Business \& Economic Statistics, 13(3).

\bibitem{gk1980} Garman, M., Klass, M. (1980). \emph{On the Estimation of
Security Price Volatilities from Historical Data}. Journal of Business, 53(1).

\bibitem{lopez2018} López de Prado, M. (2018). \emph{Advances in Financial
Machine Learning}. Wiley.

\bibitem{parkinson1980} Parkinson, M. (1980). \emph{The Extreme Value Method
for Estimating the Variance of the Rate of Return}. Journal of Business, 53(1).

\bibitem{patton2015} Patton, A., Sheppard, K. (2015). \emph{Good Volatility,
Bad Volatility: Signed Jumps and the Persistence of Volatility}. Review of
Economics and Statistics, 97(3).

\bibitem{whaley2009} Whaley, R. (2009). \emph{Understanding the VIX}. Journal
of Portfolio Management, 35(3).

\end{thebibliography}

\end{document}
